\section{Conclusion}

This study demonstrates that automated urban forest Leaf Area Index (LAI) estimation is achievable through multi-modal remote sensing (LiDAR, RGB orthophotos, near-infrared imagery) combined with machine learning regression models. Our approach addresses critical limitations of traditional manual tree inventories, which require 20-30 minutes per tree and typically exclude private property trees representing ~70\% of urban forests.

XGBoost achieved superior predictive performance (R² = 0.784, RMSE = 0.620 m²/m²), representing 39\% improvement over traditional allometric baselines. Height-related features dominated predictions (51.3\% collective importance), with logarithmic height transformation and average height as top predictors. NDVI-height interaction terms ranked prominently (6.8\% importance), indicating that spectral information modulates structural relationships beyond simple additive effects. Independent hemispherical photography validation (N=5 Ulmus trees) confirmed physical plausibility of voxelization-derived LAI estimates, with measured values (1.98-6.7 m²/m²) aligning closely with LiDAR-derived ranges (RMSE = 0.86 m²/m², ~27\% relative error comparable to field measurement uncertainty). Computational scalability analysis demonstrated city-wide feasibility: processing 25 km² (6,178 trees) required 18 hours on standard hardware (0.3 km²/hour), translating to <48 hours for full Amsterdam deployment via cloud parallelization.

This approach provides municipalities with automated, comprehensive tree inventories at significantly lower cost than manual surveys. Quantitative LAI mapping enables evidence-based tree planting strategies in heat-vulnerable neighborhoods, supporting urban heat mitigation planning as cities face intensifying thermal challenges under climate change. The methodology transfers readily to cities with standard remote sensing programs (LiDAR 25 pts/m², RGB 25cm resolution, NIR imagery).

However, several limitations warrant acknowledgment. Validation sample size remains small (N=5), constraining species-specific performance assessment. Transferability to topographically complex cities or extreme climate zones requires empirical testing. Model training focused on well-maintained municipal trees, potentially introducing bias when applied to unmaintained private trees.

Future work should expand validation across species, seasons, and climate zones; test cross-city transfer learning to quantify generalizability; and integrate LAI maps with 3D shade modeling to directly quantify cooling benefits for targeted heat mitigation interventions. Combining LAI estimates with socioeconomic data could identify shade-deficient neighborhoods for equitable greening initiatives, advancing urban climate adaptation and environmental justice objectives.

